\chapter{Adquisición analógica, visualización y registro de datos}
\section{\obj}
\capacidad
\begin{itemize}
\item Leer una entrada analógica y calcular el valor de la resistencia del termistor y la temperatura ambiente.
\item Gráficar la resistencia y la temperatura.
\item Guardar los datos en un archivo separado por comas.
\end{itemize}

\section{\mat}
\textbf{A suministrar por la Escuela:}
\begin{itemize}
\item 1 dispositivo de adquisición de datos NI myDAQ
\item 1 termistor TDC310 de \SI{10}{\kilo\ohm}
\item 1 resistor de \SI{10}{\kilo\ohm}
\item 1 capacitor de \SI{0.1}{\micro\farad}

\end{itemize}
\textbf{A suministrar por el estudiante:}
\begin{itemize}
\item 1 breadboard
\item Cables de interconexión (jumpers) macho-macho
\end{itemize}


\section{\pro}
\begin{enumerate}
\item Realice las conexiones del circuito tal y como se indica en la Figura \ref{fig:L3F1}.
\item Cree un instrumento virtual (\emph{.vi}) que realice lo siguiente:
    \begin{itemize}
        \item Leer en una estructura temporizada 500 muestras del valor de $v$ cada \SI{500}{\milli\second}, sacar el promedio de las muestras y mostrarlo en un indicador.
        \item Calcular el valor de la resistencia del termistor usando la siguiente ecuación y graficarla en tiempo real en el panel frontal.
        \begin{equation*}
            R_T = \dfrac{R}{\frac{5}{v} - 1}
        \end{equation*}
        \item Calcular el valor de la temperatura utilizando la siguiente ecuación y graficarlo en tiempo real en el panel frontal.
        \begin{equation*}
            T = \dfrac{\beta T_0}{\beta + T_0\ln{\left(\frac{R_T}{R_0}\right)}}
        \end{equation*}
        donde $\beta = 4100$, $T_0 = \SI{298.15}{\kelvin}$ y $R_0 = \SI{10}{\kilo\ohm}$
        \item Guardar los datos tomados en un archivo separado por comas (\emph{.csv}) en \emph{tiempo real}, de forma que la primera columna sea el tiempo y la segunda el valor de la resistencia del termistor y la tercera la temperatura.
        \item Generar un mensaje de alarma que indique cuando la temperatura sobrepasó los \SI{50}{\degreeCelsius}
        \item Detener la ejecución luego de \SI{60}{\second}.
    \end{itemize}
\begin{figure}[H]
    \centering
    \begin{circuitikz}
        \draw
        (5.5,3) node[ground]{}
            to[V]
        (5.5,5) node[above]{\SI{5}{\volt}}
            --
        (3,5)
            to[short]
        (3,3)
            to[sR, l_=$R_T$, v^=$v$]
        (3,0)
            to[R, l_=$R$]
        (3,-3)
            --
        (5.5,-3) node[ground]{}node[above]{AGND}
        (3,0)
            to[short]
        (1.5,0)
            to[C, l_=$C$]
        (1.5,-3)
            to[short]
        (3,-3)
        ;
        \draw
        (5.5,1.5) node[op amp, noinv input up](opamp){AI0}
        (opamp.+) --++ (-0.5,0) --++ (0,0.7) to[short,-*]++ (-0.8,0)
        (opamp.-) --++ (-0.5,0) --++ (0,-0.7) to[short,-*]++ (-0.8,0)
        ;
        \draw[dashed,blue]
        (4,7) -- (7,7)node[midway, below]{myDAQ} -- (7,-4) -- (4,-4) -- (4,7);
        \node at (-3,4) {$R_T=\SI{10}{\kilo\ohm}@\SI{298.15}{\kelvin}$};
        \node at (-3,3.5) {$R=\SI{10}{\kilo\ohm}$};
        \node at (-3,3) {$C=\SI{0.1}{\micro\farad}$};
    \end{circuitikz}
    \caption{Medición de caída de voltaje en un termistor.}
    \label{fig:L3F1}
\end{figure}

\item Incluya en su informe lo siguiente:
    \begin{itemize}
        \item Una captura de pantalla del diagrama de bloques
        \item Una captura de pantalla del panel frontal
        \item Una gráfica de $R_T$ realizada a partir del archivo \emph{.csv}, usando Python. Incluya los títulos de los ejes de forma adecuada.
        \item Una gráfica de $T$ realizada a partir del archivo \emph{.csv}, usando Python. Incluya los títulos de los ejes de forma adecuada.
    \end{itemize}

\end{enumerate}
